\section{The Finite Corner And Exact Constants}\label{sec:corner-theorem}

The purpose of this section is narrow.  It explains the exact constants behind
the measured ``corner'' in the finite numerical families, and it separates that
finite calibration from hidden localization, the desired dimension-free bound
\(H=O(\tau)\).  Throughout, \(\delta=\max_i\nu_i\) is the maximal row negative
mass, \(\tau=\sqrt\delta\), \(\rho=4\tau\), and \(\kappa=\tau/4\).  The hidden
height is
\[
  H=\max_i \distone(p_i,\conv(\W)),
\]
where \(\W\) is the set of \((\rho,\kappa)\)-well-exposed row vertices.  A top
hidden vertex is a row vertex \(v\) attaining this maximum.

\begin{definition}[Corner family and its two laws]
\modauditbadge\ \numbadge
The finite corner family is the verified d8--d12 budget-line family of exact
signed affine retractions with a hidden top vertex \(v\), a binding financier
row, and the row-witness dual certificate coming directly from the row identity
of \(v\).  This is a named verified construction class, not a parametrization
of all signed retractions; the report uses only the two laws displayed below:
\[
  H=2\delta,
  \qquad
  t^*=\frac{\nu_v}{1+\nu_v}=\frac{\delta}{1+\delta}.
\]
The first is the budget law: the apex poke is exactly twice the negativity
budget.  The second is the row-witness margin law: the row-witness certificate
is optimal at the edge and the top row has \(\nu_v=\delta\).  These laws are
family statements with a proved-mod-audit caveat; the algebraic consequences
below are exact once the laws are assumed.
\end{definition}

The corner is the finite scale at which these two family laws meet the exposure
threshold \(t^*=\kappa\).  Their intersection produces the constants below.

\begin{theorem}[Finite corner constants]
\modauditbadge\ \numbadge
On the corner family above, the edge scale is
\[
  \tau_* = 2-\sqrt 3,\qquad
  \delta_*=(2-\sqrt 3)^2.
\]
At this finite edge the wall value and the associated ratio are
\[
  \frac{H}{\tau}=2(2-\sqrt 3),
  \qquad
  \frac{\delta}{H^2}=\frac{7+4\sqrt 3}{4}.
\]
This is a finite-corner statement, not an asymptotic hidden-localization
theorem.
\end{theorem}

\begin{proof}
For the corner family, the row-witness margin law gives
\[
  t^*=\frac{\delta}{1+\delta}.
\]
Writing \(\delta=\tau^2\), hiding persists exactly while the margin is still
below the exposure threshold:
\[
  \frac{\tau^2}{1+\tau^2}<\frac{\tau}{4}.
\]
For \(\tau>0\), this is equivalent to
\[
  4\tau^2<\tau(1+\tau^2)
  \quad\Longleftrightarrow\quad
  \tau^2-4\tau+1>0.
\]
The smaller positive root is \(2-\sqrt 3\), so the finite edge is
\(\tau_* = 2-\sqrt 3\) and \(\delta_*=\tau_*^2\).  On the same family,
\(H=2\delta=2\tau^2\), hence
\[
  \frac{H}{\tau}=2\tau
  \quad\text{and at the edge}\quad
  \frac{H}{\tau}=2(2-\sqrt 3).
\]
Finally,
\[
  \frac{\delta}{H^2}
    =\frac{\delta}{4\delta^2}
    =\frac{1}{4\delta_*}
    =\frac{7+4\sqrt 3}{4}.
\]
The proof is therefore exact after the stated family laws are accepted; the
audit caveat is precisely the optimality/equality package for that family, not
the algebra above.
\end{proof}

\begin{numericalresult}[Dual-certificate confirmations]
\numbadge
The d9 dual-certificate sweep verified the same edge across cells of the
historical design parameter \(\sigma_{\mathrm{sweep}}\) (the d8 proxy for
geometric \(\sigtilde\), not the formal \(\sigoff\)):
\(0.05,0.10,0.20,0.35,0.50,0.536,0.70,1.00\), using exact idempotence
checks, multiplicity-correct \(W\), robust exposedness LPs with presolve off,
and honest \(\tau=\sqrt\delta\).  The wall cells recorded
\[
  H/\tau \approx 0.535,\qquad \delta/H^2\approx 3.484\text{--}3.497,
\]
with the same level-functional shape and the same binding financier rows from
the frame groups.

The d10 perturbation sweep found, on the probed family, the exact numerical
shadow-price law
\[
  \delta_{\min}=\frac12 g_f,\qquad g_f=H,
\]
where \(g_f\) is the canonical deficit of the binding financier row, the far
blocker that pays for the apex poke.  The d11 scale sweep then resolved the
interpretation: throughout that family,
\[
  g_f=H=2\delta
\]
over more than two decades of \(\delta\).  Thus d10 measured the budget line,
not a universal theorem about all hidden configurations.
\end{numericalresult}

The budget line explains why the corner is not a small-\(\delta\) floor.  Along
\(H=2\delta\), one has \(H/\tau=2\tau\to0\) and
\(\delta/H^2=1/(4\delta)\to\infty\) as \(\delta\to0\).
The family touches the \(0.536\) wall only because the optimizer drives the
scale up to the absolute point \(\delta_*\).  At genuinely small \(\delta\), the
same line collapses far below the old square-root picture.

\begin{downgradedstatement}[No asymptotic envelope]
\downgradedbadge
The finite corner does not prove that \(H\le 2(2-\sqrt3)\tau\) is the global
asymptotic envelope, and it does not prove the hidden-localization claim.  The
post-audit correction also retracts the earlier reading that a slack
deep-witness mass forcing (DMF) argument with \(m_*=1\) reproduces the finite
corner constants.  Such arguments lose \(O(\delta/\tau)=O(\tau)\) slack and
give only an asymptotic implication of the form \(\delta\ge aH^2\), with
constants depending on the deep-mass and slack parameters.
\end{downgradedstatement}

This distinction matters because all small-scale evidence was reframed after
the corner audit.  The old ``flat floor'' was not verified below roughly
\(\delta\approx0.06\).  The final d13 probe found outcome (c): no verified
hidden-top-vertex instance near the floor for \(\delta\le10^{-2}\), and the
best hidden-vertex heights instead track the linear law \(H\approx2\delta\).
The corner therefore calibrates a finite family; it is not evidence by itself
for a small-\(\delta\) square-root obstruction.

The robust branch variable in the later analysis is not the formal off-own-site
mass \(\sigoff\).  It is the geometric mass \(\sigtilde\): the positive
coefficient mass of \(v\) on rows outside \(C_W=\conv(\W)\), with an
\(\varepsilon\)-halo convention when a row lies near \(C_W\).  Halo rows are
automatically deep up to the same \(\varepsilon\) slack for the canonical
separator.  This convention is essential because the d12 probe used an
index-based proxy, while the final branch split is geometric.

\begin{proposition}[Small-\(\sigtilde\) height cap]
\provedbadge
If \(v\) is a top hidden vertex, \(\sigtilde\le\tau\), and \(\delta\le1/4\), then
\[
  H\le 3\tau .
\]
Consequently this branch already gives \(\delta\ge H^2/9\), and existential
DMF holds there with \(m_*=1\) and depth slack constant \(C_D=3\).
\end{proposition}

\begin{proof}
The post-design control argument gives the sharper cap
\[
  H\le 2(1+2\delta)\max(\sigtilde,\nu_v).
\]
Since \(\nu_v\le\delta=\tau^2\le\tau\) for \(\delta\le1\), the assumption
\(\sigtilde\le\tau\) gives
\[
  H\le 2(1+2\delta)\tau\le 3\tau
\]
when \(\delta\le1/4\).  The displayed quadratic bound is immediate.  With
\(C_D=3\), the DMF depth threshold \(H-C_D\delta/\tau=H-3\tau\) is nonpositive,
so every witness has mass \(1\) above that threshold.
\end{proof}

\begin{remark}[All-shallow faces and the height condition]
\modauditbadge
An exact rational \(5\times5\) construction has \(\sigtilde>0\), hidden rows,
and an optimal failed-exposedness witness supported entirely on a far shallow
row.  Its scale is \(H=1/1000\) and
\(\delta=1841/1600000\), so it lies on \(H=O(\delta)\) with \(H/\tau\) small.
Thus all-shallow optimal faces are real, but this certificate does not refute a
height-qualified DMF statement.
\end{remark}

\begin{openproblem}[Height-qualified shallow-web exclusion]
\openbadge
The remaining obstruction is to exclude hidden top vertices with
\(\sigtilde>\tau\) at height \(H>B\tau\) for a universal \(B\).  Equivalently,
one needs a signed quantitative Baake--Sumner stability statement ruling out
self-sustaining shallow hidden webs at square-root height.  The zero-negativity
anchor is understood, but no quantitative modulus is banked here.
\end{openproblem}

The pushed-witness cleanup route is also closed as a proof strategy.  Pushing
shallow witness rows through \(P\) introduces negative leakage; after
renormalization, the dual objective changes from \(B\) to \((B+N)/M_F\).  Since
\(M_F\le 1+N\), any \(N>0\) destroys optimality when \(B=t^*(v)>0\): one would
need \(M_F\ge 1+N/B\) with \(B<\kappa<1\).  The \(t^*(v)=0\) case is handled by
the exchange identity without dividing by \(B\).  The residual is therefore not
another exposedness-dual algebra trick.

\begin{warning}[Writer warning]
\downgradedbadge
Do not use the corner constants as a global theorem.  They may be cited as the
exact explanation of the finite measured corner, with the proved-mod-audit
family-law caveat above.  They may not be used to assert an asymptotic envelope,
to remove the \(\sigtilde>\tau\) web case, or to replace the open
height-qualified Baake--Sumner stability problem.
\end{warning}
